\chapter{Experimental Design}

This chapter explains how the four pricing methods are compared. It describes the contracts being priced, the neural-network configuration, the allocation of simulated paths and the measures used to assess statistical and computational efficiency.

\section{Contract specifications and monitoring profiles}

The reference contract has an initial asset price of $S_0=100$, a strike of $K=100$, a risk-free rate of $r=0.05$, a volatility of $\sigma=0.20$ and a maturity of $T=1$ year.

Six additional contracts are created by changing the strike, volatility or maturity one at a time. All other parameters remain equal to those of the reference contract.

\begin{table}[h]
\centering
\begin{adjustbox}{max width=\textwidth}
\begin{tabular}{lccc}
\hline
Contract & $K$ & $\sigma$ & $T$ \\
\hline
Reference & 100 & 0.20 & 1.0 \\
Low strike & 80 & 0.20 & 1.0 \\
High strike & 120 & 0.20 & 1.0 \\
Low volatility & 100 & 0.10 & 1.0 \\
High volatility & 100 & 0.50 & 1.0 \\
Short maturity & 100 & 0.20 & 0.5 \\
Long maturity & 100 & 0.20 & 2.0 \\
\hline
\end{tabular}
\end{adjustbox}
\end{table}

This gives a seven-contract grid in which the effect of changing each contract parameter can be examined separately. The initial asset price and risk-free rate remain fixed at $S_0=100$ and $r=0.05$ throughout.

Each contract is evaluated using 12 and 252 equally spaced monitoring dates. For the one-year reference contract, these correspond approximately to monthly and business-daily monitoring. When maturity changes, the number of dates remains fixed but the time between them changes.

The two monitoring profiles are analysed separately. A 12-date network receives 12 shocks, while a 252-date network receives 252 shocks. Networks are therefore trained and reused within each profile, rather than transferred between them.

\section{Selection of the neural-network configurations}

The neural network used for the 12-date profile has 12 inputs and 449 trainable parameters. By comparison, using the same hidden width of 32 for the 252-date profile produces 8,129 parameters. Training both networks with 5,000 simulated paths therefore provides approximately 11.14 paths per parameter for the 12-date network but only 0.62 for the 252-date network.

The initial 252-date network performed poorly when pricing compared to the GCV. To examine whether this was caused by the size of the network relative to the training sample, an analysis of different neural-network configurations was conducted. Hidden-layer widths of 8, 16 and 32 were tested using 5,000, 10,000 and 20,000 training paths.

For each configuration, checkpoints were saved at several stages during training. A checkpoint is a saved copy of the network weights and biases after a given number of epochs. Each checkpoint was evaluated using 10,000 independently simulated validation paths. The checkpoint producing the lowest average variance of
\[
f(Z)-H_\theta(Z)
\]
was selected, where $f(Z)$ is the discounted option payoff and $H_\theta(Z)$ is the network output.

Table X reports the variance-reduction ratio relative to standard Monte Carlo. A larger ratio indicates greater variance reduction.

\begin{table}[h]
\centering
\begin{adjustbox}{max width=\textwidth}
\begin{tabular}{lccc}
\hline
Hidden width & 5,000 training paths & 10,000 training paths & 20,000 training paths \\
\hline
8 & 309 & 969 & 2,393 \\
16 & 273 & 1,004 & 3,306 \\
32 & 227 & 609 & 3,510 \\
\hline
\end{tabular}
\end{adjustbox}
\end{table}

Increasing the amount of training data improved every network. With 5,000 paths, performance became worse as the hidden width increased. The width-8 network achieved a variance-reduction ratio of 309, compared with 227 for the width-32 network. All three networks also selected the relatively early checkpoint at epoch 25.

At 20,000 training paths, widths 16 and 32 performed similarly. Width 32 produced the higher average variance-reduction ratio, at 3,510 compared with 3,306. However, width 16 produced a slightly lower mean residual variance, at 0.02006 compared with 0.02086. The confidence intervals for residual variance also overlapped considerably:
\[
[0.01695,\,0.02316]
\]
for width 16 and
\[
[0.01379,\,0.02792]
\]
for width 32.

Width 16 was selected for the final 252-date experiments. It achieved similar variance reduction to width 32 while using 4,065 rather than 8,129 trainable parameters. This configuration is treated as the best configuration selected from the sensitivity analysis, rather than as a globally optimal neural network.

The final network configurations are summarised below.

\begin{table}[h]
\centering
\begin{adjustbox}{max width=\textwidth}
\begin{tabular}{lccccc}
\hline
Monitoring dates & Hidden width & Parameters & Training paths & Paths per parameter & Selected checkpoint \\
\hline
12 & 32 & 449 & 5,000 & 11.14 & 1,000 \\
252 & 16 & 4,065 & 20,000 & 4.92 & 200 \\
\hline
\end{tabular}
\end{adjustbox}
\end{table}

These configurations are used in the subsequent comparisons with standard Monte Carlo, antithetic variates and the geometric control variate.

\section{Pricing comparison}

The final neural-network configurations were compared with standard Monte Carlo, antithetic variates and the geometric control variate. Each method was applied to the seven contracts described in Section 5.1 under both monitoring profiles. The complete experiment was repeated ten times using different simulated samples.

Two comparisons were conducted. The first gave each method 50,000 final pricing observations. This allows their statistical performance to be compared using the same number of observations, but it does not account for the additional paths required to set up each method.

Standard Monte Carlo required no setup sample and used 50,000 independently simulated paths. The geometric control variate used 1,000 pilot paths to estimate its control coefficient, followed by 50,000 pricing paths. The neural control variate required either 5,000 or 20,000 training paths, depending on the monitoring profile, followed by 50,000 pricing paths.

For antithetic variates, one observation is the average payoff from a path generated using $Z$ and its paired path generated using $-Z$. Consequently, 50,000 antithetic observations require 100,000 simulated paths.

\begin{table}[h]
\centering
\begin{adjustbox}{max width=\textwidth}
\begin{tabular}{lccc}
\hline
Method & Setup paths & Pricing observations & Total path evaluations \\
\hline
Monte Carlo & 0 & 50,000 & 50,000 \\
Antithetic variates & 0 & 50,000 pairs & 100,000 \\
Geometric control variate & 1,000 & 50,000 & 51,000 \\
Neural control variate, 12 dates & 5,000 & 50,000 & 55,000 \\
Neural control variate, 252 dates & 20,000 & 50,000 & 70,000 \\
\hline
\end{tabular}
\end{adjustbox}
\end{table}

A second comparison used a fixed total budget of 50,000 simulated paths. Setup paths were included within this budget. Monte Carlo therefore retained all 50,000 paths for pricing, while antithetic variates used 25,000 pairs. After allowing for setup, the geometric control variate used 49,000 pricing paths. The 12-date neural control variate used 45,000 pricing paths, while the 252-date neural control variate used 30,000.

\begin{table}[h]
\centering
\begin{adjustbox}{max width=\textwidth}
\begin{tabular}{lccc}
\hline
Method & Setup paths & Pricing observations & Total path budget \\
\hline
Monte Carlo & 0 & 50,000 & 50,000 \\
Antithetic variates & 0 & 25,000 pairs & 50,000 \\
Geometric control variate & 1,000 & 49,000 & 50,000 \\
Neural control variate, 12 dates & 5,000 & 45,000 & 50,000 \\
Neural control variate, 252 dates & 20,000 & 30,000 & 50,000 \\
\hline
\end{tabular}
\end{adjustbox}
\end{table}

The equal-observation comparison measures how much variance each method removes once it is available for pricing. The equal-budget comparison also accounts for the simulated paths used to estimate the geometric control coefficient or train the neural network.

\section{Measuring statistical and computational performance}

The methods are first compared using their standard errors. For an estimator with observation variance $s^2$ and $N$ pricing observations, the estimated variance of the option-price estimate is
\[
\frac{s^2}{N},
\]
and its standard error is
\[
\mathrm{SE}=\frac{s}{\sqrt{N}}.
\]
A lower standard error indicates a more precise option-price estimate.

Variance reduction is reported relative to standard Monte Carlo using the variance-reduction ratio
\[
\mathrm{VRR}
=
\frac{s_{\mathrm{MC}}^2}{s_{\mathrm{method}}^2}.
\]
A ratio greater than one means that the method produces a lower observation variance than standard Monte Carlo. For example, a variance-reduction ratio of 100 means that the observation variance is 100 times smaller.

Computational performance is assessed using measured runtime. Setup time and pricing time are recorded separately. Setup includes neural-network training or estimation of the geometric control coefficient, while pricing time includes only the generation and evaluation of the final pricing observations. This distinction allows the one-off cost of setting up a method to be separated from the cost of using it for subsequent valuations.

\section{Statistical results}

Variance-reduction ratios are summarised using geometric means across the independent replications. The NCV/GCV column compares the two control variates directly. A value greater than one indicates that NCV has lower variance than GCV.

\subsection{The 12-date profile}

Table X reports the results for the 12-date network, which was trained using 5,000 paths for 1,000 epochs.

\begin{table}[h]
\centering
\begin{adjustbox}{max width=\textwidth}
\begin{tabular}{lcccc}
\hline
Contract & AV/MC & GCV/MC & NCV/MC & NCV/GCV \\
\hline
Reference & 4.16 & 1,290.2 & 22,810.8 & 17.68 \\
Low strike & 56.70 & 1,629.3 & 25,913.5 & 15.91 \\
High strike & 2.12 & 232.7 & 3,131.8 & 13.46 \\
Low volatility & 6.43 & 4,588.5 & 53,795.8 & 11.72 \\
High volatility & 3.00 & 201.4 & 3,565.3 & 17.70 \\
Short maturity & 4.05 & 2,698.6 & 30,167.7 & 11.18 \\
Long maturity & 4.29 & 598.4 & 13,949.1 & 23.31 \\
\hline
\end{tabular}
\end{adjustbox}
\end{table}

All three methods reduced variance relative to standard Monte Carlo. Antithetic variates produced the smallest improvement for most contracts, while GCV produced much larger reductions. However, NCV substantially outperformed both classical methods across the complete contract grid.

Relative to GCV, the variance reduction achieved by NCV ranged from 11.18 for the short-maturity contract to 23.31 for the long-maturity contract. NCV produced lower variance than GCV in every replication for every contract, and all confidence intervals for the NCV/GCV ratio remained above one.

For the reference contract, the standard error fell from 0.038122 under standard Monte Carlo to 0.018688 using antithetic variates and 0.001061 using GCV. NCV reduced it further to 0.000255. The 12-date results therefore provide strong evidence that a contract-specific neural control can outperform the geometric Asian control when the network is trained successfully.

\subsection{The 252-date profile}

Table X reports the results for the final 252-date network. This network had a hidden width of 16 and was trained using 20,000 paths for 200 epochs. All ten replications were completed successfully.

\begin{table}[h]
\centering
\begin{adjustbox}{max width=\textwidth}
\begin{tabular}{lccccc}
\hline
Contract & AV/MC & GCV/MC & NCV/MC & NCV/GCV & 95\% CI \\
\hline
Reference & 4.16 & 1,307.7 & 2,970.5 & 2.27 & [1.99, 2.60] \\
Low strike & 65.51 & 1,509.6 & 1,752.1 & 1.16 & [0.92, 1.47] \\
High strike & 2.11 & 201.5 & 216.8 & 1.08 & [0.91, 1.27] \\
Low volatility & 6.40 & 4,687.4 & 4,942.9 & 1.05 & [0.84, 1.32] \\
High volatility & 3.01 & 204.7 & 589.5 & 2.88 & [2.45, 3.38] \\
Short maturity & 4.04 & 2,739.5 & 4,141.1 & 1.51 & [1.32, 1.73] \\
Long maturity & 4.29 & 605.3 & 2,015.2 & 3.33 & [2.87, 3.86] \\
\hline
\end{tabular}
\end{adjustbox}
\end{table}

NCV clearly outperformed GCV for the reference, high-volatility, short-maturity and long-maturity contracts. Its largest relative advantage occurred for the long-maturity contract, where its observation variance was approximately 3.33 times lower than that of GCV.

The average NCV/GCV ratios also exceeded one for the two strike changes and the low-volatility contract. However, their confidence intervals include one, so these results do not establish a clear difference between the methods. There was no contract for which the results provided clear evidence that GCV had lower variance than NCV.

Across both monitoring profiles, the geometric control remained a strong classical benchmark. The point estimates favoured NCV for all seven 252-date contracts, although clear evidence of an advantage was obtained for only four contracts. Whether this statistical improvement justifies the additional training cost is considered in the next section.

\section{Computational performance}

The variance-reduction results do not account for the additional cost of setting up each method. GCV requires a small pilot sample to estimate its control coefficient, while NCV requires the neural network to be trained before pricing begins. This training cost is paid once, after which the fitted network can be used for repeated valuations of the same contract.

Table X summarises the measured runtimes. Standalone runtime includes both setup and pricing. Marginal pricing time measures the cost of an additional valuation after setup has been completed.

\begin{table}[h]
\centering
\begin{adjustbox}{max width=\textwidth}
\begin{tabular}{lccccc}
\hline
Monitoring dates & GCV standalone & NCV standalone & GCV marginal pricing & NCV marginal pricing & Matched-accuracy break-even \\
\hline
12 & 0.048--0.056 s & 29.5--29.9 s & 0.046--0.052 s & 0.055--0.058 s & 609--703 valuations \\
252 & 1.1--1.4 s & 37--41 s & 1.12--1.34 s & 0.80--0.98 s & 40--101 valuations \\
\hline
\end{tabular}
\end{adjustbox}
\end{table}

For the 12-date profile, training the neural network required approximately 29.5--29.9 seconds. A standalone NCV valuation was therefore considerably slower than GCV, which required less than 0.06 seconds. NCV was also slightly slower when both methods used the same number of pricing observations. Its marginal pricing time was approximately 0.055--0.058 seconds, compared with 0.046--0.052 seconds for GCV.

The computational advantage of the 12-date NCV came from its substantially lower variance. It could achieve the same standard error as GCV using fewer pricing paths. Under this matched-accuracy comparison, the initial training cost was recovered after 609--703 repeated valuations, depending on the contract.

\begin{table}[h]
\centering
\begin{adjustbox}{max width=\textwidth}
\begin{tabular}{lc}
\hline
Contract & Matched-accuracy break-even \\
\hline
Reference & 688 \\
Low strike & 684 \\
High strike & 676 \\
Low volatility & 703 \\
High volatility & 626 \\
Short maturity & 694 \\
Long maturity & 609 \\
\hline
\end{tabular}
\end{adjustbox}
\end{table}

At 1,000 repeated valuations, the projected runtime saving from using NCV rather than GCV was approximately 26\%--37\%.

The 252-date profile produced a different runtime pattern. Training the network required approximately 36--40 seconds, meaning that GCV remained much faster for a single valuation. Once training was complete, however, NCV pricing required approximately 0.80--0.98 seconds, compared with 1.12--1.34 seconds for GCV.

The 252-date NCV therefore had both a lower marginal pricing time and lower observed variance. When both methods used the same number of pricing observations, the projected break-even points across the alternative contracts ranged from 84 to 167 repeated valuations. When the methods were compared at the accuracy achieved by GCV, the range fell to 40--101 valuations.

These results show that NCV was not computationally attractive for a one-off valuation. Its training cost could only be justified when the fitted network was used repeatedly, with the required number of valuations depending on its variance reduction and marginal pricing cost.

The reported runtimes are specific to the implementation and hardware used in this study. Neural-network speed can be affected by a range of variables including the software library, batching procedure, degree of vectorisation, numerical precision and the use of CPU or GPU hardware. A more highly optimised implementation could therefore produce different training times and break-even points. However, all methods were evaluated using the same codebase, hardware and timing procedure, providing a consistent comparison within the experimental setting.

The practical relevance of these break-even points depends on how often the trained network remains suitable for pricing. If changes in market or contract parameters require a new network, the same pricing problem may not be repeated often enough to recover the training cost. This motivates the network-reuse experiment considered in the next section.

\section{Reusing the neural network}

To test whether initial training cost could be shared, the network trained for the reference contract was reused to price the six alternative contracts. Its weights and biases were held fixed, meaning that no further neural-network training was performed.

Two approaches were tested. The first reused the network with the control coefficient fixed at $\beta=1$. The second estimated a new value of $\beta$ using a 1,000-path pilot sample from the target contract. This allowed the output of the reference network to be rescaled without retraining it.

The reused network was compared with a target-specific GCV. The comparison is reported as
\[
R_{\mathrm{reuse/GCV}}
=
\frac{\mathrm{Var}(\mathrm{GCV})}
{\mathrm{Var}(\mathrm{reused\ NCV})}.
\]
A value greater than one means that the reused neural control produced lower variance than GCV. A value below one means that GCV produced lower variance.

\begin{table}[h]
\centering
\begin{adjustbox}{max width=\textwidth}
\begin{tabular}{lcccc}
\hline
Target contract & $\beta=1$, 12 dates & Estimated $\beta$, 12 dates & $\beta=1$, 252 dates & Estimated $\beta$, 252 dates \\
\hline
Low strike & 0.0031 & 0.0040 & 0.0033 & 0.0042 \\
High strike & 0.0008 & 0.0086 & 0.0007 & 0.0092 \\
Low volatility & 0.0003 & 0.0106 & 0.0003 & 0.0107 \\
High volatility & 0.0123 & 0.1267 & 0.0120 & 0.1247 \\
Short maturity & 0.0016 & 0.8410 & 0.0016 & 0.6436 \\
Long maturity & 0.0155 & 1.8304 & 0.0153 & 1.3693 \\
\hline
\end{tabular}
\end{adjustbox}
\end{table}

Reusing the network with $\beta=1$ performed poorly under both monitoring profiles. None of the fixed-coefficient controls outperformed GCV. Under the 252-date profile, fixed-coefficient reuse also performed worse than standard Monte Carlo for the high-strike contract.

Estimating a target-specific value of $\beta$ substantially improved the reused control. Relative to fixing the coefficient at one, it reduced residual variance by a geometric-mean factor of approximately 26.6 for the 12-date profile and 24.9 for the 252-date profile. Estimated-$\beta$ reuse also reduced variance relative to standard Monte Carlo for every target contract.

The estimated coefficients varied considerably across the contract grid. Under the 252-date profile, their mean values ranged from 0.207 for the high-strike contract to 2.594 for the high-volatility contract. Similar variation occurred under the 12-date profile. This explains why a coefficient fixed at one did not transfer successfully when the contract parameters changed.

Despite the improvement from estimating $\beta$, GCV remained the stronger control for five of the six target contracts under both monitoring profiles. The long-maturity contract was the only target for which reuse outperformed GCV under both monitoring profiles. Its variance advantage over GCV was 1.83 under the 12-date profile and 1.37 under the 252-date profile.

This result was consistent across all replications for the 12-date profile and all replications for the 252-date profile. Both maturity changes also transferred more successfully than the strike and volatility changes, although only two maturity changes were tested, so no general conclusion about maturity transfer can be drawn.

The results therefore show that estimating $\beta$ can correct differences in the scale of the network output, but it cannot fully adapt a network trained for one contract to a different payoff.

Direct reuse of a fixed network was not a reliable replacement for contract-specific training or GCV across the tested contract grid. However, the improvement from estimating $\beta$, together with the successful long-maturity transfer, indicates that the trained network retained some information that could be reused. Alternative approaches to adapting the network are considered in the discussion and suggestions for future research.
